%_________Appendix__________________________________

% final tutorial remarks
\ifLSRITRtutorial
	\chapter{BUSTED! Last chance to actually read the HowTo-Section!}
	
%	Make sure your thesis is well structured, that each major section does what it is supposed to do, and that the whole thing hangs together. The basic structure is often as given in this template (but other structures are possible). In particular, don't think you need to have exactly as many major sections or chapters as the list implies; sometimes it makes sense to merge things, sometimes it makes sense to move things (e.g., the literature review is in many papers deferred until after the results), sometimes it makes sense to split a logical part into several individual sections. Just use some common sense.
%	
%	Hand in your thesis at minimum \textbf{one week} before the deadline for correction. You will receive feedback for the final version and very likely have to do minor or major revisions of your writing. Plan your writing schedule to allow for these adjustments, which can have quite some impact on your grade! 
%	
%	\optional{Please have a look on our \href{https://wiki.lsr.ei.tum.de/thesiswriting_students}{thesis-guidelines} as well before submitting your \emph{final} thesis.}
	
\optional{Do not forget to check our \href{https://wiki.tum.de/display/lsritr/Students}{\underline{thesis-guidelines}} before submitting your \emph{final} thesis.}	
	
%	\section{Style and Expressions}
%	
%	Before handing in your thesis, even for an intermediate review, please perform a spellcheck and correct grammar mistakes. The report is not meant to be a narrative text. Please stick to neutral and technical style and avoid subjective or biased expressions or adjectives/adverbs such as \emph{obviously, always, very, especially well, actually, so-called etc}. Scientific writing is about precision and you should underpin your statements factually, not soften them with unnecessary qualifiers.
%	
%	... Okay enough. But please check chapter \ref{sec:Tutorial} before starting with your report. 

\paragraph{Spellcheck reminder}
Before handing in your thesis, even for an intermediate review, please perform a proper spellcheck and correct grammar mistakes. Otherwise, there is a high likelihood that your supervisor will return your report immediately without making any comments.
\fi


\chapter{Scalar Implementation of POPs}
% Insert this chapter after \appendix in the main thesis file.
\section{Scalar Implementation of the 3D-Pendulum Optimization Problem}
\label{app:controlled_3d_pendulum_scalar_pop}

This appendix expands the matrix formulation
\eqref{eq:controlled_3d_pendulum_pop} into the scalar polynomial equations used
to construct the semidefinite relaxation. Let
\begin{equation}
R_k=[r_{ij,k}]_{i,j=1}^{3},
\qquad
F_k=[f_{ij,k}]_{i,j=1}^{3},
\qquad
J_d=\operatorname{diag}(\delta_1,\delta_2,\delta_3),
\label{eq:app_3d_scalar_matrices}
\end{equation}
and define
\begin{equation}
\rho_c=
\begin{bmatrix}\rho_1&\rho_2&\rho_3\end{bmatrix}^{\top},
\qquad
\bar u_{p,k}=
\begin{bmatrix}\bar u_{p1,k}&\bar u_{p2,k}&\bar u_{p3,k}\end{bmatrix}^{\top},
\qquad
u_{p,k}=u_{\max}\bar u_{p,k}.
\label{eq:app_3d_scalar_vectors}
\end{equation}

\paragraph{Objective and Discrete Dynamics}

Writing \(R_N^{\mathrm{des}}=[r_{ij,N}^{\mathrm{des}}]\) and
\(F^{\mathrm{des}}=[f_{ij}^{\mathrm{des}}]\), the scalar objective is
\begin{align}
J
=\;&
\rho_R\sum_{i=1}^{3}\sum_{j=1}^{3}
\left(r_{ij,N}-r_{ij,N}^{\mathrm{des}}\right)^2
+\rho_F\sum_{i=1}^{3}\sum_{j=1}^{3}
\left(f_{ij,N-1}-f_{ij}^{\mathrm{des}}\right)^2
\notag\\
&+\sum_{k=0}^{N-1}\alpha_R
\sum_{i=1}^{3}\sum_{j=1}^{3}
\left(r_{ij,k}-r_{ij,N}^{\mathrm{des}}\right)^2
\notag\\
&+\sum_{k=0}^{N-2}\alpha_F
\sum_{i=1}^{3}\sum_{j=1}^{3}
\left(f_{ij,k}-f_{ij}^{\mathrm{des}}\right)^2
\notag\\
&+\sum_{k=1}^{N-1}\gamma^{-1}
\left(\bar u_{p1,k}^2+\bar u_{p2,k}^2+\bar u_{p3,k}^2\right).
\label{eq:app_3d_scalar_objective}
\end{align}

Since \(R_k^{\top}e_3=[r_{31,k}\;r_{32,k}\;r_{33,k}]^{\top}\), write the
gravity moment as
\[
\tau_{g,k}=
\begin{bmatrix}\tau_{g1,k}&\tau_{g2,k}&\tau_{g3,k}\end{bmatrix}^{\top}.
\]
Its components are
\begin{subequations}
\label{eq:app_3d_scalar_gravity}
\begin{align}
\tau_{g1,k}&=mg\left(\rho_2r_{33,k}-\rho_3r_{32,k}\right),\\
\tau_{g2,k}&=mg\left(\rho_3r_{31,k}-\rho_1r_{33,k}\right),\\
\tau_{g3,k}&=mg\left(\rho_1r_{32,k}-\rho_2r_{31,k}\right).
\end{align}
\end{subequations}
Likewise, using the thesis definition
\(u_k=R_k^{\top}e_3\times u_{p,k}\) with
\(u_{p,k}=u_{\max}\bar u_{p,k}\), the three components are
\begin{subequations}
\label{eq:app_3d_scalar_control_moment}
\begin{align}
u_{1,k}&=u_{\max}
\left(r_{32,k}\bar u_{p3,k}-r_{33,k}\bar u_{p2,k}\right),\\
u_{2,k}&=u_{\max}
\left(r_{33,k}\bar u_{p1,k}-r_{31,k}\bar u_{p3,k}\right),\\
u_{3,k}&=u_{\max}
\left(r_{31,k}\bar u_{p2,k}-r_{32,k}\bar u_{p1,k}\right).
\end{align}
\end{subequations}
After multiplying
\eqref{eq:controlled_3d_pendulum_pop_dynamics} by \(h\), the three independent
entries of the skew-symmetric dynamics constraint are
therefore
\begin{subequations}
\label{eq:app_3d_scalar_dynamics}
\begin{align}
&\delta_2f_{32,k+1}-\delta_3f_{23,k+1}
-\delta_3f_{32,k}+\delta_2f_{23,k}
-h^2\left(\tau_{g1,k+1}+u_{1,k+1}\right)=0,
\label{eq:app_3d_scalar_dynamics_1}
\\
&\delta_3f_{13,k+1}-\delta_1f_{31,k+1}
-\delta_1f_{13,k}+\delta_3f_{31,k}
-h^2\left(\tau_{g2,k+1}+u_{2,k+1}\right)=0,
\label{eq:app_3d_scalar_dynamics_2}
\\
&\delta_1f_{21,k+1}-\delta_2f_{12,k+1}
-\delta_2f_{21,k}+\delta_1f_{12,k}
-h^2\left(\tau_{g3,k+1}+u_{3,k+1}\right)=0.
\label{eq:app_3d_scalar_dynamics_3}
\end{align}
\end{subequations}
for \(k=0,\ldots,N-2\).

\paragraph{Kinematics and Rotation Constraints}

The nine scalar equations corresponding to \(R_{k+1}=R_kF_k\) are
\begin{subequations}
\label{eq:app_3d_scalar_kinematics}
\begin{align}
r_{11,k+1}-r_{11,k}f_{11,k}-r_{12,k}f_{21,k}-r_{13,k}f_{31,k}&=0,\\
r_{12,k+1}-r_{11,k}f_{12,k}-r_{12,k}f_{22,k}-r_{13,k}f_{32,k}&=0,\\
r_{13,k+1}-r_{11,k}f_{13,k}-r_{12,k}f_{23,k}-r_{13,k}f_{33,k}&=0,\\
r_{21,k+1}-r_{21,k}f_{11,k}-r_{22,k}f_{21,k}-r_{23,k}f_{31,k}&=0,\\
r_{22,k+1}-r_{21,k}f_{12,k}-r_{22,k}f_{22,k}-r_{23,k}f_{32,k}&=0,\\
r_{23,k+1}-r_{21,k}f_{13,k}-r_{22,k}f_{23,k}-r_{23,k}f_{33,k}&=0,\\
r_{31,k+1}-r_{31,k}f_{11,k}-r_{32,k}f_{21,k}-r_{33,k}f_{31,k}&=0,\\
r_{32,k+1}-r_{31,k}f_{12,k}-r_{32,k}f_{22,k}-r_{33,k}f_{32,k}&=0,\\
r_{33,k+1}-r_{31,k}f_{13,k}-r_{32,k}f_{23,k}-r_{33,k}f_{33,k}&=0.
\end{align}
\end{subequations}
for \(k=0,\ldots,N-1\).

For every \(k=0,\ldots,N\), the six independent scalar equations in
\(R_k^{\top}R_k=I_3\) are
\begin{subequations}
\label{eq:app_3d_scalar_R_orthogonality}
\begin{align}
r_{11,k}^2+r_{21,k}^2+r_{31,k}^2-1&=0,\\
r_{12,k}^2+r_{22,k}^2+r_{32,k}^2-1&=0,\\
r_{13,k}^2+r_{23,k}^2+r_{33,k}^2-1&=0,\\
r_{11,k}r_{12,k}+r_{21,k}r_{22,k}+r_{31,k}r_{32,k}&=0,\\
r_{12,k}r_{13,k}+r_{22,k}r_{23,k}+r_{32,k}r_{33,k}&=0,\\
r_{11,k}r_{13,k}+r_{21,k}r_{23,k}+r_{31,k}r_{33,k}&=0.
\end{align}
\end{subequations}
The nine right-handedness equations for \(R_k\) are
\begin{subequations}
\label{eq:app_3d_scalar_R_right_handedness}
\begin{align}
r_{21,k}r_{32,k}-r_{31,k}r_{22,k}-r_{13,k}&=0,\\
r_{31,k}r_{12,k}-r_{11,k}r_{32,k}-r_{23,k}&=0,\\
r_{11,k}r_{22,k}-r_{21,k}r_{12,k}-r_{33,k}&=0,\\
r_{22,k}r_{33,k}-r_{32,k}r_{23,k}-r_{11,k}&=0,\\
r_{32,k}r_{13,k}-r_{12,k}r_{33,k}-r_{21,k}&=0,\\
r_{12,k}r_{23,k}-r_{22,k}r_{13,k}-r_{31,k}&=0,\\
r_{23,k}r_{31,k}-r_{33,k}r_{21,k}-r_{12,k}&=0,\\
r_{33,k}r_{11,k}-r_{13,k}r_{31,k}-r_{22,k}&=0,\\
r_{13,k}r_{21,k}-r_{23,k}r_{11,k}-r_{32,k}&=0.
\end{align}
\end{subequations}

For every \(k=0,\ldots,N-1\), the corresponding six equations in
\(F_k^{\top}F_k=I_3\) are
\begin{subequations}
\label{eq:app_3d_scalar_F_orthogonality}
\begin{align}
f_{11,k}^2+f_{21,k}^2+f_{31,k}^2-1&=0,\\
f_{12,k}^2+f_{22,k}^2+f_{32,k}^2-1&=0,\\
f_{13,k}^2+f_{23,k}^2+f_{33,k}^2-1&=0,\\
f_{11,k}f_{12,k}+f_{21,k}f_{22,k}+f_{31,k}f_{32,k}&=0,\\
f_{12,k}f_{13,k}+f_{22,k}f_{23,k}+f_{32,k}f_{33,k}&=0,\\
f_{11,k}f_{13,k}+f_{21,k}f_{23,k}+f_{31,k}f_{33,k}&=0.
\end{align}
\end{subequations}
The nine right-handedness equations for \(F_k\) are
\begin{subequations}
\label{eq:app_3d_scalar_F_right_handedness}
\begin{align}
f_{21,k}f_{32,k}-f_{31,k}f_{22,k}-f_{13,k}&=0,\\
f_{31,k}f_{12,k}-f_{11,k}f_{32,k}-f_{23,k}&=0,\\
f_{11,k}f_{22,k}-f_{21,k}f_{12,k}-f_{33,k}&=0,\\
f_{22,k}f_{33,k}-f_{32,k}f_{23,k}-f_{11,k}&=0,\\
f_{32,k}f_{13,k}-f_{12,k}f_{33,k}-f_{21,k}&=0,\\
f_{12,k}f_{23,k}-f_{22,k}f_{13,k}-f_{31,k}&=0,\\
f_{23,k}f_{31,k}-f_{33,k}f_{21,k}-f_{12,k}&=0,\\
f_{33,k}f_{11,k}-f_{13,k}f_{31,k}-f_{22,k}&=0,\\
f_{13,k}f_{21,k}-f_{23,k}f_{11,k}-f_{32,k}&=0.
\end{align}
\end{subequations}

\paragraph{Bounds, Initial Values, and Normalization}

The relative-rotation and control bounds become
\begin{subequations}
\label{eq:app_3d_scalar_bounds}
\begin{align}
f_{11,k}+f_{22,k}+f_{33,k}-1
-2\cos(\Delta\theta_{\max})&\geq0,
&&k=0,\ldots,N-1,
\\
1-\bar u_{p1,k}^2-\bar u_{p2,k}^2-\bar u_{p3,k}^2&\geq0,
&&k=1,\ldots,N-1.
\end{align}
\end{subequations}
The fixed initial matrices are imposed entrywise as
\begin{equation}
r_{ij,0}=r_{ij}^{\mathrm{init}},
\qquad
f_{ij,0}=f_{ij}^{\mathrm{init}},
\qquad i,j=1,2,3.
\label{eq:app_3d_scalar_initial_values}
\end{equation}

The control scaling is already included in
\eqref{eq:app_3d_scalar_vectors},
\eqref{eq:app_3d_scalar_control_moment}, and
\eqref{eq:app_3d_scalar_bounds}. Together with
\eqref{eq:app_3d_scalar_dynamics}--
\eqref{eq:app_3d_scalar_initial_values}, these equations give the complete
scalar POP used for the controlled 3D-Pendulum relaxation.

\section{Scalar Implementation of the Acrobot Optimization Problem}
\label{app:controlled_acrobot_scalar_pop}

This appendix expands the reduced matrix formulation
\eqref{eq:controlled_acrobot_pop} into the scalar polynomial equations used to
construct the sparse semidefinite relaxation. It uses exactly the notation of
Section~\ref{sec:discrete_motion_planning_acrobot}: the normalized control and
constraint multipliers are denoted by \(\bar u_k\), \(\bar\lambda_{0,k}\), and
\(\bar\lambda_{12,k}\). The center-of-mass positions and velocities are
reconstructed from the rotations and are not optimization variables.

\paragraph{Reduced Variables and Objective}

For $i\in\{1,2\}$, let
\begin{equation}
R_{i,k}=
\begin{bmatrix}
c_{i,k}&-s_{i,k}\\
s_{i,k}& c_{i,k}
\end{bmatrix},
\qquad
F_{i,k}=
\begin{bmatrix}
a_{i,k}&-b_{i,k}\\
b_{i,k}& a_{i,k}
\end{bmatrix}.
\label{eq:app_acrobot_scalar_rotations}
\end{equation}
For the uniform-rod geometry used in the implementation,
\begin{equation}
\rho_{1,0}=\begin{bmatrix}0\\ l_1/2\end{bmatrix},
\qquad
\rho_{1,12}=\begin{bmatrix}0\\-l_1/2\end{bmatrix},
\qquad
\rho_{2,12}=\begin{bmatrix}0\\ l_2/2\end{bmatrix}.
\label{eq:app_acrobot_geometry_vectors}
\end{equation}
Hence, the holonomic geometry reconstructs the center-of-mass positions as
\begin{subequations}
\label{eq:app_acrobot_reconstructed_positions}
\begin{align}
x_{1,k}
&=p_0+
\begin{bmatrix}
\frac{l_1}{2}s_{1,k}\\[0.1em]
-\frac{l_1}{2}c_{1,k}
\end{bmatrix},
\label{eq:app_acrobot_reconstructed_x1}\\
x_{2,k}
&=p_0+
\begin{bmatrix}
l_1s_{1,k}+\frac{l_2}{2}s_{2,k}\\[0.1em]
-l_1c_{1,k}-\frac{l_2}{2}c_{2,k}
\end{bmatrix}.
\label{eq:app_acrobot_reconstructed_x2}
\end{align}
\end{subequations}
These expressions are used only for the exact elimination and for trajectory
reconstruction; $x_{1,k}$ and $x_{2,k}$ are not optimization variables.

Using \(R_{i,k+1}=R_{i,k}F_{i,k}\) and
\(R_{i,k-1}=R_{i,k}F_{i,k-1}^{\top}\), the matrix
\(A_{i,k}=F_{i,k}-2I_2+F_{i,k-1}^{\top}\), introduced in the reduced
translational dynamics of Section~\ref{sec:discrete_motion_planning_acrobot},
contains only \(F_{i,k-1}\) and
\(F_{i,k}\). Its scalar contributions are substituted directly below; no
additional hatted variables are introduced.

Let
\begin{equation}
\begin{aligned}
R_i^{\mathrm{des}}
&=\begin{bmatrix}c_i^{\mathrm{des}}&-s_i^{\mathrm{des}}\\
s_i^{\mathrm{des}}&c_i^{\mathrm{des}}\end{bmatrix},\\
F_i^{\mathrm{des}}
&=\begin{bmatrix}a_i^{\mathrm{des}}&-b_i^{\mathrm{des}}\\
b_i^{\mathrm{des}}&a_i^{\mathrm{des}}\end{bmatrix}.
\end{aligned}
\label{eq:app_acrobot_desired_rotations}
\end{equation}
Since the structured $SO(2)$ matrices satisfy
$\|R-R^{\mathrm{des}}\|_F^2=2[(c-c^{\mathrm{des}})^2+
(s-s^{\mathrm{des}})^2]$, the scalar objective is
\begin{align}
J=\;&2\rho_R\sum_{i=1}^{2}
\left[(c_{i,N}-c_i^{\mathrm{des}})^2
     +(s_{i,N}-s_i^{\mathrm{des}})^2\right]
\notag\\
&+2\rho_F\sum_{i=1}^{2}
\left[(a_{i,N-1}-a_i^{\mathrm{des}})^2
     +(b_{i,N-1}-b_i^{\mathrm{des}})^2\right]
\notag\\
&+2\alpha_R\sum_{k=0}^{N-1}\sum_{i=1}^{2}
\left[(c_{i,k}-c_i^{\mathrm{des}})^2
     +(s_{i,k}-s_i^{\mathrm{des}})^2\right]
\notag\\
&+2\alpha_F\sum_{k=0}^{N-2}\sum_{i=1}^{2}
\left[(a_{i,k}-a_i^{\mathrm{des}})^2
     +(b_{i,k}-b_i^{\mathrm{des}})^2\right]
+\gamma^{-1}\sum_{k=1}^{N-1}\bar u_k^2.
\label{eq:app_acrobot_scalar_objective}
\end{align}
There is no multiplier regularization term in the optimization problem evaluated
in this thesis.

\paragraph{Reduced Discrete Dynamics}

The normalized variables are mapped to the physical quantities exactly as in
Section~\ref{sec:discrete_motion_planning_acrobot}:
\begin{equation}
u_k^{\mathrm{phys}}=u_{\max}\bar u_k,
\qquad
\lambda_{0,k}^{(j)}
=\lambda_{\max}\bar\lambda_{0,k}^{(j)},
\qquad
\lambda_{12,k}^{(j)}
=\lambda_{\max}\bar\lambda_{12,k}^{(j)},
\quad j=1,2.
\label{eq:app_acrobot_scaling}
\end{equation}
For the geometry in \eqref{eq:app_acrobot_geometry_vectors} and the gravity
direction \(e_2=[0\;-1]^{\top}\), the four scalar components of the reduced
translational dynamics are
\begin{subequations}
\label{eq:app_acrobot_scalar_translational_dynamics}
\begin{align}
&m_1\frac{l_1}{2}
\left[
s_{1,k}(a_{1,k}+a_{1,k-1}-2)
+c_{1,k}(b_{1,k}-b_{1,k-1})
\right]
-h^2\lambda_{\max}
 \left(\bar\lambda_{0,k}^{(1)}
      +\bar\lambda_{12,k}^{(1)}\right)=0,
\label{eq:app_acrobot_translation_1x}\\
&-m_1\frac{l_1}{2}
\left[
c_{1,k}(a_{1,k}+a_{1,k-1}-2)
-s_{1,k}(b_{1,k}-b_{1,k-1})
\right]
+h^2m_1g
-h^2\lambda_{\max}
 \left(\bar\lambda_{0,k}^{(2)}
      +\bar\lambda_{12,k}^{(2)}\right)=0,
\label{eq:app_acrobot_translation_1y}\\
&m_2\Bigl[
l_1\bigl(
s_{1,k}(a_{1,k}+a_{1,k-1}-2)
+c_{1,k}(b_{1,k}-b_{1,k-1})
\bigr)
\notag\\[-0.2em]
&\hspace{5.5em}
+\frac{l_2}{2}\bigl(
s_{2,k}(a_{2,k}+a_{2,k-1}-2)
+c_{2,k}(b_{2,k}-b_{2,k-1})
\bigr)
\Bigr]
+h^2\lambda_{\max}\bar\lambda_{12,k}^{(1)}=0,
\label{eq:app_acrobot_translation_2x}\\
&-m_2\Bigl[
l_1\bigl(
c_{1,k}(a_{1,k}+a_{1,k-1}-2)
-s_{1,k}(b_{1,k}-b_{1,k-1})
\bigr)
\notag\\[-0.2em]
&\hspace{5.5em}
+\frac{l_2}{2}\bigl(
c_{2,k}(a_{2,k}+a_{2,k-1}-2)
-s_{2,k}(b_{2,k}-b_{2,k-1})
\bigr)
\Bigr]
+h^2m_2g
+h^2\lambda_{\max}\bar\lambda_{12,k}^{(2)}=0,
\label{eq:app_acrobot_translation_2y}
\end{align}
\end{subequations}
for $k=1,\ldots,N-1$.

For the geometry~\eqref{eq:app_acrobot_geometry_vectors}, the scalar moments
generated by the constraint forces are
\begin{subequations}
\label{eq:app_acrobot_scalar_constraint_moments}
\begin{align}
\mu_{1,0,k}
&=-\frac{l_1}{2}
  \lambda_{\max}\left(c_{1,k}\bar\lambda_{0,k}^{(1)}
       +s_{1,k}\bar\lambda_{0,k}^{(2)}\right),
\label{eq:app_acrobot_mu_10}\\
\mu_{1,12,k}
&=\frac{l_1}{2}
  \lambda_{\max}\left(c_{1,k}\bar\lambda_{12,k}^{(1)}
       +s_{1,k}\bar\lambda_{12,k}^{(2)}\right),
\label{eq:app_acrobot_mu_112}\\
\mu_{2,12,k}
&=-\frac{l_2}{2}
  \lambda_{\max}\left(c_{2,k}\bar\lambda_{12,k}^{(1)}
       +s_{2,k}\bar\lambda_{12,k}^{(2)}\right).
\label{eq:app_acrobot_mu_212}
\end{align}
\end{subequations}
Here,
$\mu_{1,0,k}=\rho_{1,0}\times
(R_{1,k}^{\top}\lambda_{0,k})$,
$\mu_{1,12,k}=\rho_{1,12}\times
(R_{1,k}^{\top}\lambda_{12,k})$, and
$\mu_{2,12,k}=\rho_{2,12}\times
(R_{2,k}^{\top}\lambda_{12,k})$.
Accordingly, the moment convention in the matrix POP is
$\mu_{1,k}=\mu_{1,0,k}+\mu_{1,12,k}$ and
$\mu_{2,k}=-\mu_{2,12,k}$.
The two independent scalar rotational equations are therefore
\begin{subequations}
\label{eq:app_acrobot_scalar_rotational_dynamics}
\begin{align}
&\operatorname{tr}(J_{d,1})(b_{1,k-1}-b_{1,k})
+h^2\left(
\mu_{1,0,k}+\mu_{1,12,k}-u_{\max}\bar u_k
\right)=0,
\label{eq:app_acrobot_rotation_1}\\
&\operatorname{tr}(J_{d,2})(b_{2,k-1}-b_{2,k})
+h^2\left(
u_{\max}\bar u_k-\mu_{2,12,k}
\right)=0,
\label{eq:app_acrobot_rotation_2}
\end{align}
\end{subequations}
for $k=1,\ldots,N-1$. Equations
\eqref{eq:app_acrobot_scalar_translational_dynamics} and
\eqref{eq:app_acrobot_scalar_rotational_dynamics} are all quadratic
polynomials in the reduced decision variables.

\paragraph{Kinematics, Rotation Constraints, and Bounds}

For each link, $R_{i,k+1}=R_{i,k}F_{i,k}$ gives the two independent scalar
kinematic equations
\begin{subequations}
\label{eq:app_acrobot_scalar_kinematics}
\begin{align}
c_{i,k+1}-c_{i,k}a_{i,k}+s_{i,k}b_{i,k}&=0,
\label{eq:app_acrobot_kinematics_c}\\
s_{i,k+1}-s_{i,k}a_{i,k}-c_{i,k}b_{i,k}&=0,
\label{eq:app_acrobot_kinematics_s}
\end{align}
\end{subequations}
for $i=1,2$ and $k=0,\ldots,N-1$. Owing to the structured
parameterization~\eqref{eq:app_acrobot_scalar_rotations}, the $SO(2)$
constraints reduce to
\begin{subequations}
\label{eq:app_acrobot_scalar_so2_constraints}
\begin{align}
c_{i,k}^2+s_{i,k}^2-1&=0,
&&i=1,2,\quad k=0,\ldots,N,
\label{eq:app_acrobot_so2_R}\\
a_{i,k}^2+b_{i,k}^2-1&=0,
&&i=1,2,\quad k=0,\ldots,N-1.
\label{eq:app_acrobot_so2_F}
\end{align}
\end{subequations}
The increment, control, and componentwise multiplier bounds are
\begin{subequations}
\label{eq:app_acrobot_scalar_bounds}
\begin{align}
a_{i,k}-\cos(\Delta\theta_{i,\max})&\geq0,
&&i=1,2,\quad k=1,\ldots,N-1,
\label{eq:app_acrobot_step_bounds}\\
1-\bar u_k^2&\geq0,
&&k=1,\ldots,N-1,
\label{eq:app_acrobot_control_bounds}\\
1-\bigl(\bar\lambda_{0,k}^{(j)}\bigr)^2&\geq0,
&&j=1,2,\quad k=1,\ldots,N-1,
\label{eq:app_acrobot_lambda0_bounds}\\
1-\bigl(\bar\lambda_{12,k}^{(j)}\bigr)^2&\geq0,
&&j=1,2,\quad k=1,\ldots,N-1.
\label{eq:app_acrobot_lambda12_bounds}
\end{align}
\end{subequations}
Thus, the physical quantities satisfy
$|u_k^{\mathrm{phys}}|\leq u_{\max}$ and
$|\lambda_{0,k}^{(j)}|,
 |\lambda_{12,k}^{(j)}|\leq\lambda_{\max}$.

Finally, write
\begin{equation}
R_i^{\mathrm{init}}=
\begin{bmatrix}c_i^{\mathrm{init}}&-s_i^{\mathrm{init}}\\
s_i^{\mathrm{init}}&c_i^{\mathrm{init}}\end{bmatrix},
\qquad
F_i^{\mathrm{init}}=
\begin{bmatrix}a_i^{\mathrm{init}}&-b_i^{\mathrm{init}}\\
b_i^{\mathrm{init}}&a_i^{\mathrm{init}}\end{bmatrix}.
\label{eq:app_acrobot_initial_matrices}
\end{equation}
The implementation fixes $R_{i,1}=R_i^{\mathrm{init}}$,
$F_{i,0}=F_i^{\mathrm{init}}$, and
$R_{i,0}=R_i^{\mathrm{init}}(F_i^{\mathrm{init}})^{\top}$. Entrywise, this is
\begin{subequations}
\label{eq:app_acrobot_scalar_initial_values}
\begin{align}
c_{i,1}&=c_i^{\mathrm{init}},
&s_{i,1}&=s_i^{\mathrm{init}},
&a_{i,0}&=a_i^{\mathrm{init}},
&b_{i,0}&=b_i^{\mathrm{init}},
\label{eq:app_acrobot_fixed_R1_F0}\\
c_{i,0}&=c_i^{\mathrm{init}}a_i^{\mathrm{init}}
          +s_i^{\mathrm{init}}b_i^{\mathrm{init}},
&s_{i,0}&=s_i^{\mathrm{init}}a_i^{\mathrm{init}}
          -c_i^{\mathrm{init}}b_i^{\mathrm{init}},
&&&
\label{eq:app_acrobot_reconstructed_R0}
\end{align}
\end{subequations}
for $i=1,2$. For the rest initial condition used in the numerical study,
$a_i^{\mathrm{init}}=1$ and $b_i^{\mathrm{init}}=0$, such that
$R_{i,0}=R_{i,1}=R_i^{\mathrm{init}}$. Together with the objective
\eqref{eq:app_acrobot_scalar_objective}, equations
\eqref{eq:app_acrobot_scalar_translational_dynamics}--
\eqref{eq:app_acrobot_scalar_initial_values} give the complete scalar POP
implemented for the Acrobot relaxation.

\chapter{Validation of the Reduced Acrobot Dynamics}
\label{appendix:experimental_details}

The reduction used in Section~\ref{sec:discrete_motion_planning_acrobot} is
validated against two references. The first reference is the full
maximal-coordinate LGVI used in the numerical simulation, which propagates
\(x_{i,k}\), \(v_{i,k}\), and \(R_{i,k}\) while enforcing the base and elbow
constraints. The second is the continuous joint-coordinate Acrobot model used
in Section~\ref{sec:numerical_simulation_results}, integrated with the adaptive
DOP853 method. It serves as a high-accuracy numerical approximation of the
physical dynamics; it is not a hardware measurement.

The reduced LGVI eliminates the position and velocity variables through
\eqref{eq:app_acrobot_reconstructed_positions}. It replaces the preceding
rotation by \(R_{i,k-1}=R_{i,k}F_{i,k-1}^{\top}\). Consequently, the full and reduced
LGVIs describe the same discrete equations after exact elimination. The
comparison of their reconstructed positions and absolute link rotations
confirmed agreement within the nonlinear-solver tolerance. The relevant
remaining error is therefore the common deviation of these discrete LGVI
trajectories from the continuous reference.

For the coarse discretizations \(h\in\{0.1,0.05\}\,\mathrm{s}\), a constant
elbow torque is applied over \(T_c=0.1\,\mathrm{s}\). The MPC-relevant setting
uses \(h=0.0025\,\mathrm{s}\) and \(T_c=0.025\,\mathrm{s}\). In all cases,
\[
u\in\{-20,-10,-5,0,5,10,20\}\,\mathrm{N\,m}.
\]
Let \(\Delta\theta_{i,k}\in(-\pi,\pi]\) be the shortest signed difference
between the absolute angle of link \(i\) in the LGVI and continuous
trajectories. The terminal and trajectory errors are
\begin{equation}
e_{\theta,f}=
\sqrt{\Delta\theta_{1,N_c}^{2}+\Delta\theta_{2,N_c}^{2}},
\qquad
e_{\theta,\mathrm{RMSE}}=
\sqrt{\frac{1}{2(N_c+1)}
\sum_{k=0}^{N_c}
\left(\Delta\theta_{1,k}^{2}+\Delta\theta_{2,k}^{2}\right)},
\label{eq:app_acrobot_validation_errors}
\end{equation}
where \(N_c=T_c/h\). Table~\ref{tab:app_acrobot_validation_steps} reports the
mean over all tested inputs. Since the full and reduced LGVIs coincide within
the reported precision, one common error pair is shown rather than duplicating
identical columns.

\begin{table}[ht]
\centering
\caption{Mean angular errors of the full and reduced maximal-coordinate LGVIs
relative to the continuous physical reference.}
\label{tab:app_acrobot_validation_steps}
\small
\begin{tabular}{ccccc}
\hline
\(h\,[\mathrm{s}]\) & \(T_c\,[\mathrm{s}]\) & Steps &
\(e_{\theta,f}\,[{}^\circ]\) &
\(e_{\theta,\mathrm{RMSE}}\,[{}^\circ]\)\\
\hline
0.1000 & 0.100 & 1  & 7.758 & 3.879\\
0.0500 & 0.100 & 2  & 4.202 & 1.752\\
0.0025 & 0.025 & 10 & 0.003 & 0.001\\
\hline
\end{tabular}
\end{table}

The reduction from \(h=0.1\,\mathrm{s}\) to \(h=0.05\,\mathrm{s}\) over the
same physical interval decreases both errors, although deviations of several
degrees remain for larger inputs. The last row additionally shortens the
control interval and must therefore not be interpreted as a pure step-size
convergence test. It instead measures the prediction accuracy over the interval
used by the numerical closed-loop implementation.

Table~\ref{tab:app_acrobot_validation_torque} separates the influence of the
input magnitude. For \(|u|>0\), the values are averaged over the corresponding
positive and negative torques. Each entry gives
\(e_{\theta,f}/e_{\theta,\mathrm{RMSE}}\) in degrees.

\begin{table}[ht]
\centering
\caption{Mean angular errors of the full and reduced LGVIs grouped by control
magnitude relative to the continuous reference.}
\label{tab:app_acrobot_validation_torque}
\small
\begin{tabular}{ccc}
\hline
\(|u|\,[\mathrm{N\,m}]\) &
\(h=0.05\,\mathrm{s},\ T_c=0.1\,\mathrm{s}\) &
\(h=0.0025\,\mathrm{s},\ T_c=0.025\,\mathrm{s}\)\\
\hline
0  & 0.041/0.017 & 0.000026/0.000008\\
5  & 1.860/0.773 & 0.001144/0.000346\\
10 & 4.074/1.693 & 0.002512/0.000749\\
20 & 8.752/3.658 & 0.006794/0.001953\\
\hline
\end{tabular}
\end{table}

The error increases with \(|u|\), because stronger torques produce larger
intra-step changes that are represented less accurately by a coarse finite
step. For \(h=0.0025\,\mathrm{s}\) and \(T_c=0.025\,\mathrm{s}\), the mean
terminal error remains below \(0.01^\circ\), even at
\(|u|=20\,\mathrm{N\,m}\). These results verify that the reduction preserves
the full maximal-coordinate LGVI and that the selected fine simulation scale
provides close agreement with the continuous Acrobot dynamics. The validation
concerns the prediction model itself and is independent of whether a subsequent
SDP relaxation is tight.

